\scanpage{5}

\uline{Proposition 1} Les faisceaux $\underline L^{\otimes n}$ sont coh. plats en \unsure{tte} dimension. Le faisceau $f_*(\underline L^{\otimes n})$ est localement libre, de rang $0$ si $n<0$, $1$ si $n=0$ ou $1$, $n\deleted{\unsure{+1}}$ si $n\geq1$. \deleted{\unsure{Le}} Le faisceau $\underline L$ est ample, et de façon précise $\underline L^{\otimes n}$ est \deleted{\unsure{très}} \uline{ample} \deleted{pour} \inserted{ssi} $n\geq1$, très ample ssi $n\unsure{\geq|>}3$.

La formule de RR, où $g=1$, d'où $2g-2=0$, nous montre que \uline{tout faisceau} \deleted{loc. libre} inversible sur $C$, de degré \inserted{relatif} $d\neq0$, est coh. plat, et que sa dimension est \deleted{\unsure{donc}} nulle si $d<0$, et $d\deleted{\unsure{-(g-1)}}$ si $d\geq1$. Donc la première assertion est démontrée pour $n\neq0$, et pour $n=0$ elle est connue par ailleurs, car $\underline L^{\otimes0}=\underline O_C$.
\margin{prop. préliminaire sur un faisceau inversible \unsure{quelconque}}
\inserted{On sait aussi qu'un faisceau inv. sur $C$ de degré relatif $d$ est ample ssi $d>0$, donc $\underline L^{\otimes n}$ est ample ssi $n\geq1$. \deleted{\unsure{Et puisque}}}
\deleted{Comme $\underline L$ est de degré relatif $1>0$, il est ample, \unsure{par suite} \unsure{les} $\underline L^{\otimes n}$ est très ample pour $n\unsure{\gg|>}0$.}

Si $\underline L^{\otimes n}$ est très ample, on a $n\geq1$, d'ailleurs on ne peut avoir $n=1,2$ car alors $\unsure{\mathbb P|P}(f_*(\underline L^{\otimes n}))$ est de dim relative $0$ ou $1$, donc $C$ ne peut s'y immerger (car il y serait isomorphe, absurde), donc il faut $n\geq3$.

\unsure{Ceci suffit}, car à cause de la platitude cohomologique, on est ramené au cas où $S=\operatorname{Spec}k$, $k$ un corps \inserted{alg. clos}. \deleted{Mais alors D'ailleurs, comme $\underline L$ est engendré loc. par ses sections}

Mais \deleted{alors} comme $3>(2g-2)+2=\uline{2g}=2$, un raisonnement connu n'utilisant que RR donne le résultat.

\mspage{4}
